% esferas.tex
%
% Copyright (C) 2022--2026 José A. Navarro Ramón <janr.devel@gmail.com>
% Licencia del código GPLv2
% Licencia Creative Commons Recognition Non-Commercial Share-alike.
% (CC-BY-NC-SA)

\chapter{Esferas}

En este texto veremos como representar distintos tipos de elementos sobre una esfera.

Para conseguirlo, se utiliza Lua\LaTeX , apoyado por el lenguaje \textsf{Lua}.
El propio formato del código fuente de este documento, puede servir para adaptarlo a
cualquier otra situación que necesite el usuario.

Utilizamos tres elementos importantes, que son los datos que se interpretan como
instrucciones para realizar una esfera y sus elementos concretos, el código fuente que
ejecuta esas instrucciones, y por último el programa \LaTeX\ que da la orden inicial para
que se obtenga el dibujo.   Se recomienda estudiar el código fuente

\begin{enumerate}
\item DATOS:

  Los datos contienen la información que define una esfera con sus elementos,
  como pueden ser meridianos, paralelos, segmentos de paralelo, arcos máximos, puntos,
  planos tangente y vectores en estos planos.
  
  Estos datos se encuentran como tablas \textsf{Lua}  en ficheros que, por comodidad se
  encuentran en un determinado directorio. En estos ejemplos es \texttt{'datoslua'}, pero
  podría ser cualquier otro que quiera el usuario.

  Es importante entender que se dibujarán todos los elementos que tengan imagen gráfica,
  y que se encuentren en el fichero de datos.
  
\item FUNCIONES:

  Las funciones Lua se encuentran en el fichero
  \texttt{lua/funciones\_esfera.lua}.
  
  La función principal es \texttt{lua/dibuja\_tikzesfera}, que según los datos
  que se le proporcionen, utilizará las funciones auxiliares que se necesiten.

\item Lua\LaTeX :

  Este será el lenguaje que organice el texto y las imágenes.
  Por ser un lenguaje general de confección de textos, se pueden utilizar los paquetes
  que se necesite, dentro del inmenso catálogo de ellos que hay a su disposición.

  El fichero principal de este documento es \texttt{esferas.tex}, que se encarga de
  organizar la estructura del documento.
  
  Para dibujar esferas con sus elementos, se utiliza el paquete \textsf{TikZ}, que se
  carga en el fichero de carga \texttt{esferas.pkg.sty} junto con otros que le pueden
  servir de complemento.

  Además, en el fichero de configuración y definiciones, \texttt{esferas.defs.sty},
  contiene, entre otros, el \LaTeX\ que se encargan de encontrar las funciones
  \textsf{Lua}, el comando \texttt{\textbackslash cargaDatos} que utiliza el fichero de
  datos que se le pasa como argumento, y \texttt{\textbackslash dibujaTikzEsfera}, que
  llama a la función \textsf{Lua} principal.

  Estos tres ficheros se encuentran en el directorio raíz.
\end{enumerate}

Aquí nos prepararemos para representar una esfera desprovista de elementos.

\section{Esferas planas}
La \emph{esfera plana} se representa mediante un círculo.
El campo \texttt{'esfera'} es el que se encarga de su representación.
Este consta de los siguientes valores:

\vspace{1ex}
CAMPO \texttt{'esfera'}:
\begin{itemize}
  \setlist[itemize]{noitemsep}\vspace*{-1ex}
\item \texttt{'radio'}: Radio de la circunferencia en \unit{cm} (número).
\item \texttt{'draw'}: Indica el color que se utilizará para el borde del círculo
  (cadena de texto).
\item \texttt{'lw'}: Grosor del borde de la esfera en dimensiones \LaTeX\ '\unit{pt}'
  (cadena de texto). Ejemplo: \texttt{"0.8pt"}.
\item \texttt{'fill'}: Color de relleno del círculo (cadena de texto).
  \item \texttt{'opac'}: Opacidad del relleno (número de 0 a 1).
\end{itemize}

Lo encontramos en \texttt{'esf\_bas\_01.lua'} y lo cargamos:
Lo cargamos en Lua\LaTeX\ así:

\vspace{1ex}
\begin{tabular}{|@{\hspace{-1em}}c|}\hline
  \vphantom{\rule{0.4pt}{3ex}}
  \begin{lstlisting}[basicstyle=\ttfamily\footnotesize, language=TeX]
    \cargaDatos{datoslua/esf_bas_01.lua}
  \end{lstlisting}
  \\[1ex]
  \hline
\end{tabular}

El campo o subtabla \texttt{'esfera'} en el fichero de datos pone:

\vspace{1ex}
\begin{tabular}{|@{\hspace{-1em}}c|}\hline
  \vphantom{\rule{0.4pt}{3ex}}
  \begin{lstlisting}[basicstyle=\ttfamily\footnotesize, language=TeX]
   esfera = {
      radio=2.0, draw="black!30", lw= "0.4pt", fill="black!10", opac=1.0,
   },
  \end{lstlisting}
  \\[1ex]
  \hline
\end{tabular}

\vspace{1ex}
El código para  generar la imagen de la esfera podría ser el siguiente, aunque se
pueden utilizar otras variantes de Lua\LaTeX :

\vspace{1ex}
\begin{tabular}{|@{\hspace{-1em}}c|}\hline
  \vphantom{\rule{0.4pt}{3ex}}
  \begin{lstlisting}[basicstyle=\ttfamily\footnotesize, language=TeX]
    \cargaDatos{datoslua/esf_bas_01.lua}
    \begin{figure}[ht]
      \begin{center}
        \begin{tikzpicture}[scale=1.0]
          \dibujaTikzEsfera
        \end{tikzpicture}
      \end{center}
      \caption{Esfera básica plana, con contorno y relleno.} 
    \end{figure}
  \end{lstlisting}
  \\[1ex]
  \hline
\end{tabular}

El resultado se aprecia en la figura \ref{fig:esf-bas_plana1}

\begin{figure}[ht]
  \centering
  \cargaDatos{datoslua/esf_bas_01.lua}
  \begin{tikzpicture}[scale=1.0,]
    \dibujaTikzEsfera
  \end{tikzpicture}
  \caption{Esfera básica plana, con contorno y relleno.}
  \label{fig:esf-bas_plana1}
\end{figure}

 
Si desactivamos \texttt{fill = "--"}, conseguiremos una circunferencia sin relleno,
como se comprueba en el fichero \texttt{esf\_bas\_02.lua}:

\vspace{1ex}
\begin{tabular}{|@{\hspace{-1em}}c|}\hline
  \vphantom{\rule{0.4pt}{3ex}}
  \begin{lstlisting}[basicstyle=\ttfamily\footnotesize, language=TeX]
   esfera = {radio=2.0, draw="black!40", lw="0.8pt", fill="--" , opac=1.0,},    
  \end{lstlisting}
  \\[1ex]
  \hline
\end{tabular}

Pero, si desactivamos \texttt{draw = "--"}, como en \texttt{esf\_bas\_03.lua}, tendríamos
una esfera sin contorno, con solo relleno:

\vspace{1ex}
\begin{tabular}{|@{\hspace{-1em}}c|}\hline
  \vphantom{\rule{0.4pt}{3ex}}
  \begin{lstlisting}[basicstyle=\ttfamily\footnotesize, language=TeX]
   esfera = {radio=2.0, draw="--", lw="0.8pt", fill="black!10" , opac=1.0,},    
  \end{lstlisting}
  \\[1ex]
  \hline
\end{tabular}

\vspace{1ex}
Ambas esferas se pueden ver reunidas en la figura \ref{fig:esf-bas_plana23}:
 
\begin{figure}[ht]
  \centering
  \begin{minipage}{0.45\textwidth}
    \cargaDatos{datoslua/esf_bas_02.lua}
    \begin{tikzpicture}[scale=1.0,]
      \dibujaTikzEsfera
    \end{tikzpicture}
  \end{minipage}
  \hfill
  \begin{minipage}{0.45\textwidth}
    \cargaDatos{datoslua/esf_bas_03.lua}
    \centering
    \begin{tikzpicture}[scale=1.0,]
      \dibujaTikzEsfera
    \end{tikzpicture}
  \end{minipage}
  \caption{A la izquierda, esfera sin relleno. A la derecha sin contorno.}
  \label{fig:esf-bas_plana23}
\end{figure}


\section{Esferas con sombreado}
Para dotar a la esfera de perspectiva, se puede añadir un cierto efecto de sombreado a
través del campo \texttt{'smbresfera'}, que complementa al campo \texttt{'esfera'}, que
define el valor \texttt{'radio'} necesario para el dibujo de la esfera.
Este campo tiene los siguientes valores:

\vspace{1ex}
CAMPO \texttt{'smbresfera'}:
\begin{itemize}
  \setlist[itemize]{noitemsep}\vspace*{-1ex}
\item \texttt{'ballcolor'}: Indica el color del reflejo en la esfera.
  (cadena de texto).
  \item \texttt{'opac'}: Opacidad del relleno (número de 0 a 1).
\end{itemize}

En el fichero de datos \texttt{esf\_bas\_02.lua}, los campos \texttt{'esfera'} y \texttt{'smbresfera'} tienen:

\vspace{1ex}
\begin{tabular}{|@{\hspace{-1em}}c|}\hline
  \vphantom{\rule{0.4pt}{3ex}}
  \begin{lstlisting}[basicstyle=\ttfamily\footnotesize, language=TeX]
   esfera = {radio=2.0, draw="white", lw="0pt", fill="white", opac=0,},
   smbresfera = {ballcolor="green!50", opac=0.35},
  \end{lstlisting}
  \\[1ex]
  \hline
\end{tabular}

\vspace{1ex}
El código para  generar la imagen de la esfera podría ser el siguiente, aunque se
pueden utilizar otras variantes de Lua\LaTeX :

\vspace{1ex}
\begin{tabular}{|@{\hspace{-1em}}c|}\hline
  \vphantom{\rule{0.4pt}{3ex}}
  \begin{lstlisting}[basicstyle=\ttfamily\footnotesize, language=TeX]
    \cargaDatos{datoslua/esf_bas_02.lua}
    \begin{figure}[ht]
      \begin{center}
        \begin{tikzpicture}[scale=1.0]
          \dibujaTikzEsfera
        \end{tikzpicture}
      \end{center}
      \caption{Esfera básica plana.} 
    \end{figure}
  \end{lstlisting}
  \\[1ex]
  \hline
\end{tabular}

\vspace{1ex}
Como resultado,
%figura \ref{fig:esf-sombr_verde},
se genera una superficie esférica
limpia, pues el campo \texttt{'esfera'} es el único en el fichero de datos.

\vspace{1ex}
Si cambiamos los valores de los campos, como en \texttt{esf\_bas\_03.lua}:

\vspace{1ex}
\begin{tabular}{|@{\hspace{-1em}}c|}\hline
  \vphantom{\rule{0.4pt}{3ex}}
  \begin{lstlisting}[basicstyle=\ttfamily\footnotesize, language=TeX]
   esfera = {radio=2.0, draw="white", lw= "0pt", fill="white", opac=0,},
   smbresfera = {ballcolor="white", opac=0.35},
  \end{lstlisting}
  \\[1ex]
  \hline
\end{tabular}


\begin{figure}[ht]
  \centering
  \begin{minipage}{0.45\textwidth}
    \cargaDatos{datoslua/esf_smbr_01.lua}
    \begin{tikzpicture}[scale=1.0,]
      \dibujaTikzEsfera
    \end{tikzpicture}
  \end{minipage}
  \hfill
  \begin{minipage}{0.45\textwidth}
    \cargaDatos{datoslua/esf_smbr_02.lua}
    \centering
    \begin{tikzpicture}[scale=1.0,]
      \dibujaTikzEsfera
    \end{tikzpicture}
  \end{minipage}
  \caption{Esferas sombreadas que muestran una cierta perspectiva, en verde y en gris.}
  \label{fig:esf-sombreadas}
\end{figure}








%%% Local Variables:
%%% coding: utf-8
%%% mode: latex
%%% TeX-engine: luatex
%%% TeX-master: "../elementos_esfera.tex"
%%% End:


